% =============================================================================
% abstract.tex — Abstract + Highlights (v0.3, fixed for standard article class)
% FIX: \begin{highlights} is elsarticle-specific; replaced by a plain section
%      + itemize so this compiles under the standard article class.
% =============================================================================

\begin{abstract}
Biochar is among the most scalable engineered carbon dioxide removal (CDR)
pathways, but its economics depend on carbon accounting rather than energy
sales, and existing models treat carbon policy as exogenous. We present a
multi-product biochar supply chain model for Wisconsin (72 county nodes, 13
feedstocks, drying plus slow pyrolysis at 300/500~$^\circ$C, price-responsive
demand) in the market-clearing framework of Sampat et al.~\cite{sampat2019},
where mass-balance duals are spatial equilibrium prices. We embed three policy
architectures: baseline-and-credit, cap-and-trade with an endogenous carbon
price, and a tiered tax with Verra VM0044 removal crediting. The no-policy
equilibrium yields 2.30~Mt biochar/yr, \$708~M/yr surplus, and 1.60~Mt
CO$_2$e/yr net removal. A linear carbon tax leaves this optimum unchanged and
net emission caps never bind, because the system is demand-constrained and
net-negative; removal credits instead activate the \$250/t bulk market at
\$25/tCO$_2$e and mobilize nearly all near-term biomass at \$200/tCO$_2$e
(surplus \$3.36~B/yr). The cap's endogenous marginal abatement cost
(\$351--599/tCO$_2$e) exceeds EU ETS and voluntary credit prices, and the
VM0044 H/C gate reshapes temperature selection. Policy design, not stringency,
determines the scale, technology, and spatial structure of biochar CDR.
\end{abstract}

\paragraph{Keywords} biochar supply chain; carbon dioxide removal; market
clearing; inherent values; endogenous carbon price; baseline-and-credit;
cap-and-trade; tiered carbon tax; VM0044; mixed-integer programming

\section*{Highlights}
\begin{itemize}
\item Show that a linear carbon tax leaves the demand-constrained optimum
unchanged.
\item Trace an endogenous abatement cost of \$351--599/tCO$_2$e, above market
prices.
\item Show that net emission caps never bind for a net-negative system.
\item Find that removal credits activate the bulk market at \$25/tCO$_2$e.
\item Show the VM0044 H/C gate converts the system to 500$^\circ$C pyrolysis.
\end{itemize}
